This section provides a list of all proof rules supported by
Alethe.  To make this long list more accessible, the section
first lists multiple classes of proof rules.  The classes are not
mutually exclusive: for example, the \proofRule{la_generic} rule is
both a linear arithmetic rule and introduces a tautology.
The number in brackets is the position of the rule in the overall list
of proof rules.
Table~\ref{rule-tab:special} lists rules that serve a special purpose.
Table~\ref{rule-tab:tautologies} lists all rules that introduce
tautologies.  That is, regular rules that do not use premises.

The subsequent section, starting at \pageref{sec:alethe:rules-list}, defines
all rules and provides example proofs for complicated rules.
The index of proof rules at \pageref{sec:alethe:rules-index} can be used
to quickly find the definition of rules.

\subsection{Classifications of the Rules}
\label{sec:alethe:rules-overview}

\begin{xltabular}{\linewidth}{l X}
\caption{Special rules.}
\label{rule-tab:special}\\
	Rule & Description \\
	\hline
	\ruleref{assume}   & Repetition of an input assumption. \\
	\ruleref{hole}     & Placeholder for rules not defined here. \\
	\ruleref{subproof} & Concludes a subproof and discharges local assumptions. \\
\end{xltabular}

\begin{xltabular}{\linewidth}{l X}
\caption{Resolution and related rules.}
\label{rule-tab:resolution}\\
	Rule & Description \\
	\hline
	\ruleref{resolution} & Chain resolution of two or more clauses. \\
	\ruleref{th_resolution} & As \proofRule{resolution}, but used by theory solvers. \\
	\ruleref{tautology} & Simplification of tautological clauses to $\top$. \\
	\ruleref{contraction} & Removal of duplicated literals. \\
\end{xltabular}

\begin{xltabular}{\linewidth}{l X}
\caption{Rules introducing tautologies.}
\label{rule-tab:tautologies}\\
	Rule & Description \\
	\hline
\ruleref{true} & $\top$ \\
\ruleref{false} & $\neg\bot$ \\
\ruleref{not_not} & $\neg(\neg\neg\varphi), \varphi$ \\
\ruleref{la_generic} & Tautologous disjunction of linear inequalities. \\
\ruleref{lia_generic} & Tautologous disjunction of linear integer inequalities. \\
\ruleref{la_disequality} & $t_1 ≈ t_2 \lor \neg (t_1 \leq t_2) \lor \neg (t_2 \leq t_1)$ \\
\ruleref{la_totality} & $t_1 \leq t_2 \lor t_2 \leq t_1$ \\
\ruleref{la_tautology} & A trivial linear tautology. \\
\ruleref{forall_inst} & Quantifier instantiation. \\
\ruleref{refl} & Reflexivity after applying the context. \\
\ruleref{eq_reflexive} & $t ≈ t$ without context. \\
\ruleref{eq_transitive} & $\neg (t_1 ≈ t_2) , \dots , \neg (t_{n-1} ≈ t_n) , t_1 ≈ t_n$ \\
\ruleref{eq_congruent} & $\neg (t_1 ≈ u_1) , \dots , \neg (t_n ≈ u_n) , f(t_1, \dots, t_n) ≈ f(u_1, \dots, u_n)$ \\
\ruleref{eq_congruent_pred} & $\neg (t_1 ≈ u_1) , \dots , \neg (t_n ≈ u_n) , P(t_1, \dots, t_n) ≈ P(u_1, \dots, u_n)$ \\
\ruleref{qnt_cnf} & Clausification of a quantified formula. \\
\ruleref{and_pos} & $\neg (\varphi_1 \land \dots \land \varphi_n) , \varphi_k $ \\
\ruleref{and_neg} & $ (\varphi_1 \land \dots \land \varphi_n), \neg\varphi_1  , \dots , \neg\varphi_n $ \\
\ruleref{or_pos} & $ \neg (\varphi_1 \lor \dots \lor \varphi_n) , \varphi_1 , \dots , \varphi_n $ \\
\ruleref{or_neg} & $ (\varphi_1 \lor \dots \lor \varphi_n) , \neg \varphi_k $ \\
\ruleref{xor_pos1} & $ \neg (\varphi_1 \,\lsymb{xor}\, \varphi_2) , \varphi_1 , \varphi_2 $ \\
\ruleref{xor_pos2} & $ \neg (\varphi_1 \,\lsymb{xor}\, \varphi_2), \neg \varphi_1, \neg \varphi_2 $ \\
\ruleref{xor_neg1} & $ \varphi_1 \,\lsymb{xor}\, \varphi_2, \varphi_1 , \neg \varphi_2 $ \\
\ruleref{xor_neg2} & $ \varphi_1 \,\lsymb{xor}\, \varphi_2, \neg \varphi_1 , \varphi_2 $ \\
\ruleref{implies_pos} & $ \neg (\varphi_1 \rightarrow \varphi_2), \neg \varphi_1 , \varphi_2 $ \\
\ruleref{implies_neg1} & $ \varphi_1 \rightarrow \varphi_2, \varphi_1 $ \\
\ruleref{implies_neg2} & $ \varphi_1 \rightarrow \varphi_2, \neg \varphi_2 $ \\
\ruleref{equiv_pos1} & $\neg (\varphi_1 ≈ \varphi_2) , \varphi_1 , \neg \varphi_2$ \\
\ruleref{equiv_pos2} & $\neg (\varphi_1 ≈ \varphi_2) , \neg \varphi_1 , \varphi_2$ \\
\ruleref{equiv_neg1} & $\varphi_1 ≈ \varphi_2 , \neg \varphi_1 , \neg \varphi_2$ \\
\ruleref{equiv_neg2} & $\varphi_1 ≈ \varphi_2 , \varphi_1 , \varphi_2$ \\
\ruleref{ite_pos1} & $\neg (\lsymb{ite} \varphi_1\;\varphi_2\;\varphi_3) , \varphi_1 , \varphi_3$ \\
\ruleref{ite_pos2} & $\neg (\lsymb{ite} \varphi_1\;\varphi_2\;\varphi_3) , \neg \varphi_1 , \varphi_2$ \\
\ruleref{ite_neg1} & $\lsymb{ite} \varphi_1\;\varphi_2\;\varphi_3 , \varphi_1 , \neg \varphi_3$ \\
\ruleref{ite_neg2} & $\lsymb{ite} \varphi_1\;\varphi_2\;\varphi_3 , \neg \varphi_1 , \neg \varphi_2$ \\
\ruleref{connective_def} & Definition of the Boolean connectives. \\
\ruleref{and_simplify} & Simplification of a conjunction. \\
\ruleref{or_simplify} & Simplification of a disjunction. \\
\ruleref{not_simplify} & Simplification of a Boolean negation. \\
\ruleref{implies_simplify} & Simplification of an implication. \\
\ruleref{equiv_simplify} & Simplification of an equivalence. \\
\ruleref{bool_simplify} & Simplification of combined Boolean connectives. \\
\ruleref{ac_simp} & Flattening of nested $\lor$ or $\land$. \\
\ruleref{ite_simplify} & Simplification of if-then-else. \\
\ruleref{qnt_simplify} & Simplification of constant quantified formulas. \\
\ruleref{qnt_join} & Joining of consecutive quantifiers. \\
\ruleref{qnt_rm_unused} & Removal of unused quantified variables. \\
\ruleref{eq_simplify} & Simplification of equality. \\
\ruleref{div_simplify} & Simplification of division. \\
\ruleref{prod_simplify} & Simplification of products. \\
\ruleref{unary_minus_simplify} & Simplification of the unary minus. \\
\ruleref{minus_simplify} & Simplification of subtractions. \\
\ruleref{sum_simplify} & Simplification of sums. \\
\ruleref{comp_simplify} & Simplification of arithmetic comparisons. \\
\ruleref{distinct_elim} & Elimination of the distinction predicate. \\
\ruleref{la_rw_eq} & $(t ≈ u) ≈ (t \leq u \land u \leq t)$ \\
\ruleref{nary_elim} & Replace $n$-ary operators with binary application. \\
\end{xltabular}

\begin{xltabular}{\linewidth}{l X}
\caption{Linear arithmetic rules.}
\label{rule-tab:la-tauts}\\
	Rule & Description \\
	\hline
\ruleref{la_generic} & Tautologous disjunction of linear inequalities. \\
\ruleref{lia_generic} & Tautologous disjunction of linear integer inequalities. \\
\ruleref{la_disequality} & $t_1 ≈ t_2 \lor \neg (t_1 \leq t_2) \lor \neg (t_2 \leq t_1)$ \\
\ruleref{la_totality} & $t_1 \leq t_2 \lor t_2 \leq t_1$ \\
\ruleref{la_tautology} & A trivial linear tautology. \\
\ruleref{la_rw_eq} & $(t ≈ u) ≈ (t \leq u \land u \leq t)$ \\
\ruleref{div_simplify} & Simplification of division. \\
\ruleref{prod_simplify} & Simplification of products. \\
\ruleref{unary_minus_simplify} & Simplification of the unary minus. \\
\ruleref{minus_simplify} & Simplification of subtractions. \\
\ruleref{sum_simplify} & Simplification of sums. \\
\ruleref{comp_simplify} & Simplification of arithmetic comparisons. \\
\end{xltabular}

\begin{xltabular}{\linewidth}{l X}
\caption{Quantifier handling.}
\label{rule-tab:quants}\\
	Rule & Description \\
	\hline
\ruleref{forall_inst} & Instantiation of a universal variable. \\
\ruleref{bind} & Renaming of bound variables. \\
\ruleref{sko_ex} & Skolemization of existential variables. \\
\ruleref{sko_forall} & Skolemization of universal variables. \\
\ruleref{qnt_cnf} & Clausification of quantified formulas. \\
\ruleref{qnt_simplify} & Simplification of constant quantified formulas. \\
\ruleref{onepoint} & The one-point rule. \\
\ruleref{qnt_join} & Joining of consecutive quantifiers. \\
\ruleref{qnt_rm_unused} & Removal of unused quantified variables. \\
\end{xltabular}

\begin{xltabular}{\linewidth}{l X}
\caption{Skolemization rules.}
\label{rule-tab:skos}\\
	Rule & Description \\
	\hline
\ruleref{sko_ex} & Skolemization of existential variables. \\
\ruleref{sko_forall} & Skolemization of universal variables. \\
\end{xltabular}

\begin{xltabular}{\linewidth}{l X}
\caption{Clausification rules.  These rules can be used to perform propositional
clausification.}
\label{rule-tab:clausification}\\
	Rule & Description \\
	\hline
\ruleref{and} & And elimination. \\
\ruleref{not_or} & Elimination of a negated disjunction. \\
\ruleref{or} & Disjunction to clause. \\
\ruleref{not_and} & Distribution of negation over a conjunction. \\
\ruleref{xor1} & From $\,\lsymb{xor}\, \varphi_1 \varphi_2$ to $\varphi_1, \varphi_2$. \\
\ruleref{xor2} & From $\,\lsymb{xor}\, \varphi_1 \varphi_2$ to $\neg\varphi_1, \neg\varphi_2$. \\
\ruleref{not_xor1} & From $\neg(\,\lsymb{xor}\, \varphi_1 \varphi_2)$ to $\varphi_1, \neg\varphi_2$. \\
\ruleref{not_xor2} & From $\neg(\,\lsymb{xor}\, \varphi_1 \varphi_2)$ to $\neg\varphi_1, \varphi_2$. \\
\ruleref{implies} & From $ \varphi_1\rightarrow\varphi_2 $ to $\neg\varphi_1 , \varphi_2 $. \\
\ruleref{not_implies1} & From $\neg (\varphi_1\rightarrow\varphi_2)$ to $\varphi_1$. \\
\ruleref{not_implies2} & From $\neg (\varphi_1\rightarrow\varphi_2)$ to $\neg\varphi_2$. \\
\ruleref{equiv1} & From $ \varphi_1≈\varphi_2$ to $\neg\varphi_1 , \varphi_2$. \\
\ruleref{equiv2} & From $ \varphi_1≈\varphi_2$ to $\varphi_1 , \neg\varphi_2$. \\
\ruleref{not_equiv1} & From $\neg(\varphi_1≈\varphi_2)$ to $\varphi_1 , \varphi_2$. \\
\ruleref{not_equiv2} & From $\neg(\varphi_1≈\varphi_2)$ to $\neg\varphi_1 , \neg\varphi_2$. \\
\ruleref{and_pos} & $\neg (\varphi_1 \land \dots \land \varphi_n), \varphi_k$\\
\ruleref{and_neg} & $(\varphi_1 \land \dots \land \varphi_n), \neg\varphi_1, \dots , \neg\varphi_n $ \\
\ruleref{or_pos} & $\neg (\varphi_1 \lor \dots \lor \varphi_n) , \varphi_1 , \dots
  , \varphi_n $ \\
\ruleref{or_neg} &
$ (\varphi_1 \lor \dots \lor \varphi_n) , \neg \varphi_k $\\
\ruleref{xor_pos1} &
$ \neg (\varphi_1 \,\lsymb{xor}\, \varphi_2) , \varphi_1 , \varphi_2 $ \\
\ruleref{xor_pos2} &
$ \neg (\varphi_1 \,\lsymb{xor}\, \varphi_2)
, \neg \varphi_1, \neg \varphi_2 $ \\
\ruleref{xor_neg1} &
$ \varphi_1 \,\lsymb{xor}\, \varphi_2, \varphi_1 , \neg \varphi_2 $ \\
\ruleref{xor_neg2} &
$ \varphi_1 \,\lsymb{xor}\, \varphi_2, \neg \varphi_1 , \varphi_2 $ \\
\ruleref{implies_pos} &
$ \neg (\varphi_1 \rightarrow \varphi_2), \neg \varphi_1 , \varphi_2 $ \\
\ruleref{implies_neg1} &
$ \varphi_1 \rightarrow \varphi_2, \varphi_1 $ \\
\ruleref{implies_neg2} &
$ \varphi_1 \rightarrow \varphi_2, \neg \varphi_2 $ \\
\ruleref{equiv_pos1} &
$\neg (\varphi_1 ≈ \varphi_2) , \varphi_1 , \neg \varphi_2$ \\
\ruleref{equiv_pos2} &
$\neg (\varphi_1 ≈ \varphi_2) , \neg \varphi_1 , \varphi_2$ \\
\ruleref{equiv_neg1} &
$\varphi_1 ≈ \varphi_2 , \neg \varphi_1 , \neg \varphi_2$ \\
\ruleref{equiv_neg2} &
$\varphi_1 ≈ \varphi_2 , \varphi_1 , \varphi_2$ \\
\ruleref{let} & Elimination of the $\lsymb{let}$ operator.  \\
\ruleref{distinct_elim} & Elimination of the $\lsymb{distinct}$ operator.  \\
\ruleref{nary_elim} & Elimination of n-ary application of operators. \\
\end{xltabular}

\begin{xltabular}{\linewidth}{l X}
\caption{Simplification rules. These rules represent typical operator-level
simplifications.}
\label{rule-tab:simplification}\\
	Rule & Description \\
	\hline
\ruleref{connective_def} & Definition of the Boolean connectives. \\
\ruleref{and_simplify} & Simplification of a conjunction. \\
\ruleref{or_simplify} & Simplification of a disjunction. \\
\ruleref{not_simplify} & Simplification of a Boolean negation. \\
\ruleref{implies_simplify} & Simplification of an implication. \\
\ruleref{equiv_simplify} & Simplification of an equivalence. \\
\ruleref{bool_simplify} & Simplification of combined Boolean connectives. \\
\ruleref{ac_simp} & Simplification of nested disjunctions and conjunctions. \\
\ruleref{ite_simplify} & Simplification of if-then-else. \\
\ruleref{qnt_simplify} & Simplification of constant quantified formulas. \\
\ruleref{onepoint} & The one-point rule. \\
\ruleref{qnt_join} & Joining of consecutive quantifiers. \\
\ruleref{qnt_rm_unused} & Removal of unused quantified variables. \\
\ruleref{eq_simplify} & Simplification of equalities. \\
\ruleref{div_simplify} & Simplification of division. \\
\ruleref{prod_simplify} & Simplification of products. \\
\ruleref{unary_minus_simplify} & Simplification of the unary minus. \\
\ruleref{minus_simplify} & Simplification of subtractions. \\
\ruleref{sum_simplify} & Simplification of sums. \\
\ruleref{comp_simplify} & Simplification of arithmetic comparisons. \\
\ruleref{qnt_simplify} & Simplification of constant quantified formulas. \\
\end{xltabular}

\subject{Rule List}
\label{sec:alethe:rules-list}

\begin{RuleDescription}{assume}
\begin{AletheX}
$i$. & \ctxsep & $\varphi$ & \currule \\
\end{AletheX}
	where $\varphi$ corresponds up to the orientation of equalities
	to a formula asserted in the input problem.
	\ruleparagraph{Remark.}
	This rule can not be used by the
	\inlineAlethe{(step }\dots\inlineAlethe{)} command. Instead it corresponds to the dedicated
	\inlineAlethe{(assume }\dots\inlineAlethe{)} command.
\end{RuleDescription}

\endinput

\startRuleDescription{hole}
\startAlethe
\setupTABLE[c][3][width=15.3em]
\bTR\bTD $i$.\eTD\bTD \ctxsep  \eTD\bTD
$\varphi$    \eTD\bTD (\currule\; p_1, \dots, p_n)\, [a_1, \dots, a_n] \eTD\eTR
\stopAlethe

\noindent
where $\varphi$ is any well-formed formula.

This rule can be used to express holes in the proof.  It can be used by
solvers as a placeholder for proof steps that are not yet expressed
by the proof rules in this document.  A proof checker {\em must not}
accept a proof as valid that contains this rule even if the checker can
somehow check this rule.  However, it is possible for checkers to have
a dedicated status for proofs that contain this rule and are otherwise
valid.  Any other tool can accept or reject proofs that contain this rule.

The premises and arguments are arbitrary, but must follow the syntax
for premises and arguments.

\stopRuleDescription

\startRuleDescription{true}
\startAlethe
\setupTABLE[c][3][width=27.5em]
\bTR\bTD $i$.\eTD\bTD \ctxsep  \eTD\bTD
$\top$    \eTD\bTD \currule \eTD\eTR
\stopAlethe
\stopRuleDescription

\startRuleDescription{false}
\startAlethe
\setupTABLE[c][3][width=27.5em]
\bTR\bTD $i$.\eTD\bTD \ctxsep  \eTD\bTD
$\bot$    \eTD\bTD \currule \eTD\eTR
\stopAlethe
\stopRuleDescription

\startRuleDescription{not_not}

\startAlethe
\setupTABLE[c][3][width=27.5em]
\bTR\bTD $i$.\eTD\bTD \ctxsep  \eTD\bTD
$\neg (\neg \neg \varphi)$, $\varphi$ \eTD\bTD \currule \eTD\eTR
\stopAlethe

\tightpara{Remark.} This rule is useful to remove double negations from a
clause by resolving a clause with the double negation on $\varphi$.
\stopRuleDescription

\page
\startRuleDescription{th_resolution}
This rule is the resolution of two or more clauses.
\startAlethe
\setupTABLE[c][3][width=20.1em]
\bTR\bTD $i_1$.\eTD\bTD \ctxsep \eTD\bTD $l^1_1$, \dots, $l^1_{k^1}$ \eTD\bTD ($\ruleType{...}$) \eTD\eTR
\aletheLine
\bTR\bTD $i_n$.\eTD\bTD \ctxsep \eTD\bTD $l^n_1$, \dots, $l^n_{k^n}$ \eTD\bTD $ (\ruleType{...}$) \eTD\eTR
\bTR\bTD $j$.\eTD\bTD   \ctxsep \eTD\bTD $l^{r_1}_{s_1}$, \dots, $l^{r_m}_{s_m}$ \eTD\bTD ($\currule\;i_1,\dots, i_n)$ \eTD\eTR 
\stopAlethe

\noindent
where $l^{r_1}_{s_1} , \dots , l^{r_m}_{s_m}$ are from $l^{i}_{j}$ and
are the result of a chain of predicate resolution steps on the clauses of
$i_1$ to $i_n$. It is possible that $m = 0$, i.e. that
the result is the empty clause.  When performing resolution steps, the
rule implicitly merges repeated negations at the start of the formulas
$l^{i}_{j}$.  For example, the formulas $\neg\neg\neg P$ and $\neg\neg P$
can serve as pivots during resolution.  The first formula is interpreted as
$\neg P$ and the second as just $P$ for the purpose of performing resolution
steps.

This rule is only used when the resolution step is not emitted by the SAT solver.
See the equivalent \proofRule{resolution} rule for the rule emitted by the
SAT solver.

\tightpara{Remark.} The definition given here is very general.  The motivation
for this to ensure the definition covers all possible resolution steps generated
by the existing proof generation code in veriT.  It will be restricted after
a full review of the code.  As a consequence of this checking this rule is
theoretically NP-complete.  In practice, however, the \currule-steps
produced by {\verit} are simple. Experience with reconstructing the step in
{\isabelle} shows that checking can done by naive decision procedures. The
vast majority of \currule-steps are binary resolution steps.
\stopRuleDescription

\startRuleDescription{resolution}
This rule is equivalent to the \proofRule{th_resolution} rule, but it is
emitted by the SAT solver instead of theory reasoners. The differentiation
serves only informational purpose.

\stopRuleDescription

\startRuleDescription{tautology}

\startAlethe
\setupTABLE[c][3][width=25em]
\bTR\bTD $i$.\eTD\bTD \ctxsep  \eTD\bTD
$l_1$, \dots, $l_k$, \dots, $l_m$, \dots, $l_n$ \eTD\bTD (\dots) \eTD\eTR
\bTR\bTD $j$.\eTD\bTD \ctxsep  \eTD\bTD
$\top$ \eTD\bTD (\currule\; $i$) \eTD\eTR
\stopAlethe

\noindent
  and $l_k$, $l_m$ are such that
  \startformula\startalign
  \NC l_k \NC= \underbrace{\neg \dots \neg}_o \varphi \NR
  \NC l_m \NC= \underbrace{\neg \dots \neg}_p \varphi
  \stopalign\stopformula
  and one of $o, p$\/ is odd and the other even.  Either can be $0$.
\stopRuleDescription

\startRuleDescription{contraction}
\startAlethe
\setupTABLE[c][3][width=24.4em]
\bTR\bTD $i$.\eTD\bTD \ctxsep  \eTD\bTD
$l_1$, \dots, $l_n$ \eTD\bTD (\dots) \eTD\eTR
\bTR\bTD $j$.\eTD\bTD \ctxsep  \eTD\bTD
$l_{k_1}$, \dots, $l_{k_m}$ \eTD\bTD (\currule\; $i$) \eTD\eTR
\stopAlethe

\noindent
  where $m \leq n$ and  $k_1 \dots k_m$ is a monotonic map to $1 \dots n$
  such that $l_{k_1} \dots l_{k_m}$ are pairwise distinct and
  $\{l_1, \dots, l_n\} = \{l_{k_1}, \dots, l_{k_m}\}$.
  Hence, this rule removes duplicated literals.
\stopRuleDescription


\startRuleDescription{subproof}
The \currule{} rule completes a subproof and discharges local
assumptions. Every subproof starts with some \proofRule{assume} steps. The
last step of the subproof is the conclusion.

\startAlethe
\bTR\bTD $i_1$.\eTD\bTD[leftframe=on] \ctxsep \eTD\bTD $\varphi_1$ \eTD\bTD $\proofRule{assume}$ \eTD\eTR
\aletheLineB
\bTR\bTD $i_n$.\eTD\bTD[leftframe=on] \ctxsep \eTD\bTD $\varphi_n$ \eTD\bTD $\proofRule{assume}$ \eTD\eTR 
\aletheLineB
\bTR[bottomframe=on]\bTD[bottomframe=off] $j$.\eTD\bTD[leftframe=on] \ctxsep \eTD\bTD $\psi$ \eTD\bTD (\dots) \eTD\eTR 
\bTR\bTD k.\eTD\bTD \ctxsep \eTD\bTD $\neg\varphi_1$,\dots, $\neg\varphi_n$, $\psi$  \eTD\bTD $\proofRule{subproof}$ \eTD\eTR 
\stopAlethe

\stopRuleDescription

\startRuleDescription{la_generic}
A step of the \currule{} rule represents a tautological clause of linear
disequalities.  It can be checked by showing that the conjunction of
the negated disequalities is unsatisfiable. After the application of
some strengthening rules, the resulting conjunction is unsatisfiable,
even if integer variables are assumed to be real variables.

A linear inequality is of term of the form
\startformula
\sum_{i=0}^{n}c_i\times{}t_i +
d_1\bowtie \sum_{i=n+1}^{m} c_i\times{}t_i + d_2
\stopformula
where $\bowtie\;\in \{≈, <,
>, \leq, \geq\}$, where $m\geq n$, $c_i, d_1, d_2$ are either integer or real
constants, and for each $i$ $c_i$ and $t_i$ have the same sort. We will write
$s_1 \bowtie s_2$.

Let $l_1,\dots, l_n$ be linear inequalities and $a_1, \dots, a_n$
rational numbers, then a {\currule} step has the form

\startAlethe
\setupTABLE[c][3][width=21em]
\bTR\bTD $i$.\eTD\bTD \ctxsep    \eTD\bTD
$\varphi_1$, \dots , $\varphi_o$ \eTD\bTD \currule\, [$a_1$, \dots, $a_o$] \eTD\eTR
\stopAlethe

\noindent
where $\varphi_i$ is either $\neg l_i$ or $l_i$, but never $s_1
≈ s_2$.

If the current theory does not have rational numbers, then the $a_i$ are
printed using integer division. They should, nevertheless, be interpreted
as rational numbers. If $d_1$ or $d_2$ are $0$, they might not be printed.

To check the unsatisfiability of the negation of $\varphi_1\lor \dots
\lor \varphi_o$ one performs the following steps for each literal. For
each $i$, let $\varphi := \varphi_i$ and $a := a_i$.

\startitemize[n,packed]
    \item If $\varphi = s_1 > s_2$, then let $\varphi := s_1 \leq s_2$.
    	If $\varphi = s_1 \geq s_2$, then let $\varphi := s_1 < s_2$.
    	If $\varphi = s_1 < s_2$, then let $\varphi := s_1 \geq s_2$. 
    	If $\varphi = s_1 \leq s_2$, then let $\varphi := s_1 > s_2$. This negates
    	the literal.
   	\item If $\varphi = \neg (s_1 \bowtie s_2)$, then let $\varphi := s_1 \bowtie s_2$. 
    \item Replace $\varphi$ by $\sum_{i=0}^{n}c_i\times{}t_i - \sum_{i=n+1}^{m} c_i\times{}t_i
    \bowtie d$ where $d := d_2 - d_1$.
    \item Now $\varphi$ has the form $s_1 \bowtie d$. If all
    variables in $s_1$ are integer sorted: replace $\bowtie d$ according to
    \in{Table}[tbl:lageneric].
    \item If $\bowtie$ is $≈$ replace $l$ by
    $\sum_{i=0}^{m}a\times{}c_i\times{}t_i ≈ a\times{}d$, otherwise replace it by
    $\sum_{i=0}^{m}|a|\times{}c_i\times{}t_i ≈ |a|\times{}d$.
\stopitemize

\placetable[here][tbl:lageneric]{Strengthening rules for \currule{}.}{
	\bTABLE
	\setupTABLE[r][each][frame=off]
	\setupTABLE[c][each][roffset=0.5em]
	\setupTABLE[r][1][bottomframe=on]
\bTR\bTD $\bowtie$ \eTD\bTD If $d$ is an integer \eTD\bTD Otherwise \eTD\eTR
\bTR\bTD $>$       \eTD\bTD $\geq d + 1$ \eTD\bTD $\geq \lfloor d\rfloor + 1$  \eTD\eTR
\bTR\bTD $\geq$    \eTD\bTD $\geq d$     \eTD\bTD $\geq \lfloor d\rfloor + 1$  \eTD\eTR
	\eTABLE
}

Finally, the sum of the resulting literals is trivially contradictory.
The sum
\startformula
    \sum_{k=1}^{o}\sum_{i=1}^{m^o}c_i^k*t_i^k \bowtie \sum_{k=1}^{o}d^k
\stopformula
where $c_i^k$ is the constant $c_i$ of literal $l_k$, $t_i^k$ is the term
$t_i$ of $l_k$, and $d^k$ is the constant $d$ of $l_k$. The operator
$\bowtie$ is $≈$ if all operators are $≈$, $>$ if all are
either $≈$ or $>$, and $\geq$ otherwise. The $a_i$ must be such
that the sum on the left-hand side is $0$ and the right-hand side is $>0$ (or
$\geq 0$ if $\bowtie$ is $>$).
\stopRuleDescription

% For some reason verbatim doesn't work in the enumerate environment
\startruleexample
A simple \currule{} step in the logic LRA might look like this:

\startSMTCode
(step t10 (cl (not (> (f a) (f b))) (not (= (f a) (f b))))
    :rule la_generic :args (1.0 (- 1.0)))
\stopSMTCode

To verify this we have to check the insatisfiability of $(f\,a) > (f\,b) \land
(f\,a) ≈ (f\,b)$ (step 2). After step~3 we get $(f\,a) - (f\,b) > 0 \land
(f\,a) - (f\,b) ≈ 0$. Since we don't have an integer sort in this logic step~4 does
not apply. Finally, after step~5 the conjunction is $(f\,a) - (f\,b) > 0 \land
- (f\,a) + (f\,b) ≈ 0$. This sums to $0 > 0$, which is a contradiction.
\stopruleexample

\startruleexample
The following \currule{} step is from a {\ss QF\_UFLIA} problem:
\startSMTCode
(step t11 (cl (not (<= f3 0)) (<= (+ 1 (* 4 f3)) 1))
    :rule la_generic :args (1 (div 1 4)))
\stopSMTCode
After normalization we get $-f_3 \geq 0 \land 4\times f_3 > 0$.
This time step~4 applies and we can strengthen this to
$-f_3 \geq 0 \land 4\times f_3 \geq 1$ and after multiplication we get
$-f_3 \geq 0 \land f_3 \geq ¼$. Which sums to the contradiction
$¼ \geq 0$.
\stopruleexample

\startRuleDescription{lia_generic}
This rule is a placeholder rule for integer arithmetic solving. It takes the
same form as \proofRule{la_generic}, without the additional arguments.

\startAlethe
\setupTABLE[c][3][width=24em]
\bTR\bTD $i$.\eTD\bTD \ctxsep    \eTD\bTD
$\varphi_1$, \dots , $\varphi_o$ \eTD\bTD (\currule) \eTD\eTR
\stopAlethe

\noindent
with $\varphi_i$ being linear inequalities. The disjunction
$\varphi_1\lor \dots \lor \varphi_n$ is a tautology in the theory of linear
integer arithmetic.

\tightpara{Remark.} Since this rule can introduce a disjunction of arbitrary
linear integer inequalities without any additional hints, proof checking
can be NP-hard. Hence, this rule should be avoided when possible.

\stopRuleDescription

\startRuleDescription{la_disequality}
\startAlethe
\setupTABLE[c][3][width=24em]
\bTR\bTD $i$.\eTD\bTD \ctxsep    \eTD\bTD
$t_1 ≈ t_2 \lor \neg (t_1 \leq t_2) \lor \neg (t_2 \leq t_1)$
\eTD\bTD (\currule) \eTD\eTR
\stopAlethe
\stopRuleDescription

\startRuleDescription{la_totality}
\startAlethe
\setupTABLE[c][3][width=24.5em]
\bTR\bTD $i$.\eTD\bTD \ctxsep    \eTD\bTD
$t_1 \leq t_2 \lor t_2 \leq t_1$
\eTD\bTD (\currule) \eTD\eTR
\stopAlethe
\stopRuleDescription

\startRuleDescription{la_tautology}
This rule is a linear arithmetic tautology which can be checked without
sophisticated reasoning. It has either the form

\startAlethe
\setupTABLE[c][3][width=24.5em]
\bTR\bTD $i$.\eTD\bTD \ctxsep    \eTD\bTD
$\varphi$
\eTD\bTD (\currule) \eTD\eTR
\stopAlethe

\noindent
where $\varphi$ is either a linear inequality $s_1 \bowtie s_2$
or $\neg(s_1 \bowtie s_2)$. After performing step 1 to 3 of the process for
checking the \proofRule{la_generic} the result is trivially unsatisfiable.

The second form handles bounds on linear combinations. It is binary clause:
\startAlethe
\setupTABLE[c][3][width=24.5em]
\bTR\bTD $i$.\eTD\bTD \ctxsep    \eTD\bTD
$\varphi_1 \lor \varphi_2$ % Checked: this is a proper disjunction not a cl
\eTD\bTD (\currule) \eTD\eTR
\stopAlethe

It can be checked by using the procedure for \proofRule{la_generic} with
while setting the arguments to $1$. Informally, the rule follows one of several
cases:
\startitemize[packed]
    \item $\neg (s_1 \leq d_1) \lor s_1 \leq d_2$ where $d_1 \leq d_2$
    \item $s_1 \leq d_1 \lor \neg (s_1 \leq d_2)$ where $d_1 = d_2$
    \item $\neg (s_1 \geq d_1) \lor s_1 \geq d_2$ where $d_1 \geq d_2$
    \item $s_1 \geq d_1 \lor \neg (s_1 \geq d_2)$ where $d_1 = d_2$
    \item $\neg (s_1 \leq d_1) \lor \neg(s_1 \geq d_2)$ where $d_1 < d_2$
\stopitemize
The inequalities $s_1 \bowtie d$ are are the result of applying normalization
as for the rule \proofRule{la_generic}.
\stopRuleDescription

\startRuleDescription{bind}
  The \currule{} rule is used to rename bound variables.

\startAletheS
\setupTABLE[c][4][width=15.3em]
\aletheLineS
\bTR[bottomframe=on] \bTD[bottomframe=off] $j$.\eTD\bTD[leftframe=on]
  $\Gamma, y_1,\dots, y_n,  x_1 \mapsto y_1, \dots ,  x_n \mapsto y_n$
	\eTD\bTD \ctxsep
  \eTD\bTD
  $\varphi ≈ \varphi'$
  \eTD\bTD (\dots) \eTD\eTR
\bTR\bTD $k$.\eTD \bTD \eTD
    \bTD \ctxsep \eTD\bTD
    $\forall x_1, \dots, x_n.\varphi ≈ \forall y_1, \dots, y_n. \varphi'$
    \eTD\bTD \currule{} \eTD\eTR
     \eTD\bTD \eTD\eTR
\stopAletheS

  \noindent
  where the variables $y_1, \dots, y_n$ is not free in $\forall x_1,
  \dots, x_n.\varphi$.
\stopRuleDescription

\page
\startRuleDescription{sko_ex}
The \currule{} rule skolemizes existential quantifiers.

\startAletheS
\setupTABLE[c][4][width=17em]
\aletheLineS
\bTR[bottomframe=on] \bTD[bottomframe=off] $j$.\eTD\bTD[leftframe=on]
$\Gamma, x_1 \mapsto \varepsilon_1, \dots ,  x_n \mapsto \varepsilon_n$
	\eTD\bTD \ctxsep
  \eTD\bTD
  $\varphi ≈ \psi$
  \eTD\bTD (\dots) \eTD\eTR
\bTR\bTD $k$.\eTD \bTD \eTD
    \bTD \ctxsep \eTD\bTD
    $\exists x_1, \dots, x_n.\varphi ≈ \psi$
    \eTD\bTD \currule{} \eTD\eTR
     \eTD\bTD \eTD\eTR
\stopAletheS

\noindent
where $\varepsilon_i$ stands for $\varepsilon x_i. (\exists x_{i+1}, \dots,
x_n. \varphi[x_1\mapsto \varepsilon_1,\dots,x_{i-1}\mapsto\varepsilon_{i-1}])$.
\stopRuleDescription

\startRuleDescription{sko_forall}
The \currule{} rule skolemizes universal quantifiers.

\startAletheS
\setupTABLE[c][4][width=11em]
\aletheLineS
\bTR[bottomframe=on] \bTD[bottomframe=off] $j$.\eTD\bTD[leftframe=on]
$\Gamma, x_1 \mapsto (\varepsilon x_1.\neg\varphi), \dots,  x_n \mapsto (\varepsilon x_n.\neg\varphi)$
	\eTD\bTD \ctxsep
  \eTD\bTD
  $\varphi ≈ \psi$
  \eTD\bTD (\dots) \eTD\eTR
\bTR\bTD $k$.\eTD \bTD \eTD
    \bTD \ctxsep \eTD\bTD
$\forall x_1, \dots, x_n.\varphi ≈ \psi$
    \eTD\bTD \currule{} \eTD\eTR
     \eTD\bTD \eTD\eTR
\stopAletheS

\stopRuleDescription


\startRuleDescription{forall_inst}
\startAlethe
\setupTABLE[c][3][width=18em]
\setupTABLE[c][4][width=11em,align=flushright]
\bTR\bTD $i$.\eTD\bTD \ctxsep    \eTD\bTD
$\neg (\forall x_1, \dots, x_n. P) \lor P[x_1\mapsto t_1]\dots[x_n\mapsto t_n]$
\eTD\bTD \currule\, [$(x_{k_1}, t_{k_1})$, \dots, $(x_{k_n}, t_{k_n})$] \eTD\eTR
\stopAlethe

\noindent
where $k_1, \dots, k_n$ is a permutation of $1, \dots, n$ and $x_i$ and
$k_i$ have the same sort. The arguments $(x_{k_i}, t_{k_i})$ are printed as
\inlineAlethe{(:= xki tki)}.

\tightpara{Remark.} In \in{Section}[sec:qsimp-proof] we discuss an alternative
quantifier instantiation rule.
The resulting proof would be more fine grained and this
would also be an opportunity to provide a proof for the clausification
as currently done by \proofRule{qnt_cnf}.
\stopRuleDescription

\startRuleDescription{refl}
\startAletheS
\setupTABLE[c][4][width=28em]
\bTR\bTD $j$.\eTD\bTD \ctxsep \eTD\bTD $\Gamma$ \eTD\bTD
$t_1 ≈ t_2 $
\eTD\bTD \currule \eTD\eTR
\stopAletheS

\noindent
where, if $\sigma = \subst(\Gamma)$,
the terms $t_1\sigma$ and $t_2$ are
syntactically equal up to the orientation of equalities.

\tightpara{Remark.} This is the only rule that requires the application of
the context.
\stopRuleDescription

\page
\startRuleDescription{trans}
\startAletheS
\setupTABLE[c][4][width=22em]
\bTR\bTD $i_1$.\eTD\bTD \ctxsep \eTD\bTD $\Gamma$ \eTD\bTD
$t_1 ≈ t_2 $
\eTD\bTD (\dots) \eTD\eTR
\bTR\bTD $i_2$.\eTD\bTD \ctxsep \eTD\bTD $\Gamma$ \eTD\bTD
$t_2 ≈ t_3 $
\eTD\bTD (\dots) \eTD\eTR
\aletheLineSB
\bTR\bTD $i_n$.\eTD\bTD \ctxsep \eTD\bTD $\Gamma$ \eTD\bTD
$t_n ≈ t_{n+1} $
\eTD\bTD (\dots) \eTD\eTR
\bTR\bTD $j$.\eTD\bTD \ctxsep \eTD\bTD $\Gamma$ \eTD\bTD
$t_1 ≈ t_{n+1}$
\eTD\bTD (\currule\; $i_1$, \dots, $i_n$) \eTD\eTR
\stopAletheS

\stopRuleDescription

\startRuleDescription{cong}
\startAletheS
\setupTABLE[c][4][width=22em]
\bTR\bTD $i_1$.\eTD\bTD \ctxsep \eTD\bTD $\Gamma$ \eTD\bTD
$t_1 ≈ u_1 $
\eTD\bTD (\dots) \eTD\eTR
\bTR\bTD $i_2$.\eTD\bTD \ctxsep \eTD\bTD $\Gamma$ \eTD\bTD
$t_2 ≈ u_2 $
\eTD\bTD (\dots) \eTD\eTR
\aletheLineSB
\bTR\bTD $i_n$.\eTD\bTD \ctxsep \eTD\bTD $\Gamma$ \eTD\bTD
$t_n ≈ u_n $
\eTD\bTD (\dots) \eTD\eTR
\bTR\bTD $j$.\eTD\bTD \ctxsep \eTD\bTD $\Gamma$ \eTD\bTD
$(f\,t_1\,\cdots\,t_n) ≈ (f\,u_1\,\cdots\,u_n)$
\eTD\bTD (\currule\; $i_1$, \dots, $i_n$) \eTD\eTR
\stopAletheS

\noindent
where $f$\, is any $n$-ary function symbol of appropriate sort.
\stopRuleDescription

\startRuleDescription{eq_reflexive}
\startAlethe
\setupTABLE[c][3][width=24.5em]
\bTR\bTD $i$.\eTD\bTD \ctxsep    \eTD\bTD
$t ≈ t$
\eTD\bTD \currule \eTD\eTR
\stopAlethe
\stopRuleDescription

\startRuleDescription{eq_transitive}
\startAlethe
\setupTABLE[c][3][width=24.5em]
\bTR\bTD $i$.\eTD\bTD \ctxsep    \eTD\bTD
$\neg (t_1 ≈ t_2)$, \dots , $\neg (t_{n-1} ≈ t_n)$,
$t_1 ≈ t_n$
\eTD\bTD \currule \eTD\eTR
\stopAlethe
\stopRuleDescription

\startRuleDescription{eq_congruent}
\startAlethe
\setupTABLE[c][3][width=24.5em]
\bTR\bTD $i$.\eTD\bTD \ctxsep    \eTD\bTD
$\neg (t_1 ≈ u_1)$, \dots ,$\neg (t_n ≈ u_n)$,
$(f\,t_1\, \cdots\,t_n) ≈ (f\,u_1\, \cdots\,u_n)$
\eTD\bTD \currule \eTD\eTR
\stopAlethe
\stopRuleDescription

\startRuleDescription{eq_congruent_pred}
\startAlethe
\setupTABLE[c][3][width=22em]
\bTR\bTD $i$.\eTD\bTD \ctxsep    \eTD\bTD
$\neg (t_1 ≈ u_1)$, \dots ,$\neg (t_n ≈ u_n)$,
$(P\,t_1\, \cdots\,t_n) ≈ (P\,u_1\, \cdots\,u_n)$
\eTD\bTD \currule \eTD\eTR
\stopAlethe

\noindent
where $P$\, is a function symbol with co-domain sort $\lsymb{Bool}$.
\stopRuleDescription

\page
\startRuleDescription{qnt_cnf}
\startAlethe
\setupTABLE[c][3][width=27em]
\bTR\bTD $i$.\eTD\bTD \ctxsep    \eTD\bTD
$\neg(\forall x_1, \dots, x_n. \varphi) \lor \forall x_{k_1}, \dots, x_{k_m}.\varphi'$
\eTD\bTD \currule \eTD\eTR
\stopAlethe

This rule expresses clausification of a term under a universal
quantifier. This is used by conflicting instantiation. $\varphi'$ is one of the clause
of the clause normal form of $\varphi$. The variables $x_{k_1}, \dots, x_{k_m}$ are
a permutation of $x_1, \dots, x_n$ plus additional variables added by prenexing
$\varphi$. Normalization is performed in two phases. First, the negative normal form
is formed, then the result is prenexed. The result of the first step is $\Phi(\varphi, 1)$
where:

\startformula\startalign
\NC  \Phi(\neg \varphi, 1) \NC:= \Phi(\varphi, 0)\NR
\NC  \Phi(\neg \varphi, 0) \NC:= \Phi(\varphi, 1)\NR
\NC  \Phi(\varphi_1 \lor\dots\lor\varphi_n, 1) \NC:= \Phi(\varphi_1, 1)\lor\dots\lor\Phi(\varphi_n, 1)\NR
\NC  \Phi(\varphi_1 \land\dots\land\varphi_n, 1) \NC:= \Phi(\varphi_1, 1)\land\dots\land\Phi(\varphi_n, 1)\NR
\NC  \Phi(\varphi_1 \lor\dots\lor\varphi_n, 0) \NC:= \Phi(\varphi_1, 0)\land\dots\land\Phi(\varphi_n, 0)\NR
\NC  \Phi(\varphi_1 \land\dots\land\varphi_n, 0) \NC:= \Phi(\varphi_1, 0)\lor\dots\lor\Phi(\varphi_n, 0)\NR
\NC  \Phi(\varphi_1 \rightarrow \varphi_2, 1) \NC:= (\Phi(\varphi_1, 0) \lor \Phi(\varphi_2, 1)) \land
                                                    (\Phi(\varphi_2, 0) \lor \Phi(\varphi_1, 1))\NR
\NC  \Phi(\varphi_1 \rightarrow \varphi_2, 0) \NC:= (\Phi(\varphi_1, 1) \land \Phi(\varphi_2, 0)) \lor
                                                    (\Phi(\varphi_2, 1) \land \Phi(\varphi_1, 0))\NR
\NC  \Phi(\lsymb{ite}\,\varphi_1\,\varphi_2\,\varphi_3, 1) \NC:=
             (\Phi(\varphi_1, 0) \lor \Phi(\varphi_2, 1)) \land (\Phi(\varphi_1, 1) \lor \Phi(\varphi_3, 1))\NR
\NC  \Phi(\lsymb{ite}\,\varphi_1\,\varphi_2\,\varphi_3, 0) \NC:=
             (\Phi(\varphi_1, 1) \land \Phi(\varphi_2, 0)) \lor (\Phi(\varphi_1, 0) \land \Phi(\varphi_3, 0))\NR
\NC  \Phi(\forall x_1, \dots, x_n. \varphi, 1) \NC:= \forall x_1, \dots, x_n. \Phi(\varphi, 1)\NR
\NC  \Phi(\exists x_1, \dots, x_n. \varphi, 1) \NC:= \exists x_1, \dots, x_n. \Phi(\varphi, 1)\NR
\NC  \Phi(\forall x_1, \dots, x_n. \varphi, 0) \NC:= \exists x_1, \dots, x_n. \Phi(\varphi, 0)\NR
\NC  \Phi(\exists x_1, \dots, x_n. \varphi, 0) \NC:= \forall x_1, \dots, x_n. \Phi(\varphi, 0)\NR
\NC  \Phi(\varphi, 1) \NC:= \varphi\NR
\NC  \Phi(\varphi, 0) \NC:= \neg\varphi
\stopalign\stopformula

\tightpara{Remark.} This is a placeholder rule that combines the many steps
done during clausification.

\stopRuleDescription

\page
\startRuleDescription{and}
\startAlethe
\setupTABLE[c][3][width=28em]
\bTR\bTD $i$.\eTD\bTD \ctxsep    \eTD\bTD
$\varphi_1 \land \cdots \land \varphi_n$
\eTD\bTD (\dots) \eTD\eTR
\bTR\bTD $j$.\eTD\bTD \ctxsep    \eTD\bTD
$\varphi_k$
\eTD\bTD (\currule\;$i$) \eTD\eTR
\stopAlethe

\noindent
and $1 \leq k \leq n$.
\stopRuleDescription

\startRuleDescription{not_or}
\startAlethe
\setupTABLE[c][3][width=27em]
\bTR\bTD $i$.\eTD\bTD \ctxsep    \eTD\bTD
$\neg (\varphi_1 \lor \cdots \lor \varphi_n)$
\eTD\bTD (\dots) \eTD\eTR
\bTR\bTD $j$.\eTD\bTD \ctxsep    \eTD\bTD
$\neg \varphi_k$
\eTD\bTD (\currule\;$i$) \eTD\eTR
\stopAlethe

\noindent
and $1 \leq k \leq n$.
\stopRuleDescription

\startRuleDescription{or}
\startAlethe
\setupTABLE[c][3][width=28em]
\bTR\bTD $i$.\eTD\bTD \ctxsep    \eTD\bTD
$\varphi_1 \lor \cdots \lor \varphi_n$
\eTD\bTD (\dots) \eTD\eTR
\bTR\bTD $j$.\eTD\bTD \ctxsep    \eTD\bTD
$\varphi_1, \dots, \varphi_n$
\eTD\bTD (\currule\;$i$) \eTD\eTR
\stopAlethe

\tightpara{Remark.} This rule deconstructs the \inlineAlethe{or} operator
into a clause denoted by \inlineAlethe{cl}.
\stopRuleDescription

\startruleexample
An application of the \currule{} rule.

\startSMTCode
(step t15 (cl (or (= a b) (not (<= a b)) (not (<= b a))))
    :rule la_disequality)
(step t16 (cl     (= a b) (not (<= a b)) (not (<= b a)))
    :rule or :premises (t15))
\stopSMTCode
\stopruleexample


\startRuleDescription{not_and}
\startAlethe
\setupTABLE[c][3][width=26em]
\bTR\bTD $i$.\eTD\bTD \ctxsep    \eTD\bTD
$\neg (\varphi_1 \land \dots \land \varphi_n)$
\eTD\bTD (\dots) \eTD\eTR
\bTR\bTD $j$.\eTD\bTD \ctxsep    \eTD\bTD
$\neg\varphi_1 , \dots , \neg\varphi_n$
\eTD\bTD (\currule\;$i$) \eTD\eTR
\stopAlethe

\stopRuleDescription

\startRuleDescription{xor1}
\startAlethe
\setupTABLE[c][3][width=28em]
\bTR\bTD $i$.\eTD\bTD \ctxsep    \eTD\bTD
$\lsymb{xor}\,\varphi_1\,\varphi_2$
\eTD\bTD (\dots) \eTD\eTR
\bTR\bTD $j$.\eTD\bTD \ctxsep    \eTD\bTD
$\varphi_1, \varphi_2$
\eTD\bTD (\currule\;$i$) \eTD\eTR
\stopAlethe
\stopRuleDescription

\page
\startRuleDescription{xor2}
\startAlethe
\setupTABLE[c][3][width=28em]
\bTR\bTD $i$.\eTD\bTD \ctxsep    \eTD\bTD
$\lsymb{xor}\,\varphi_1\,\varphi_2$
\eTD\bTD (\dots) \eTD\eTR
\bTR\bTD $j$.\eTD\bTD \ctxsep    \eTD\bTD
$\neg\varphi_1 , \neg\varphi_2$
\eTD\bTD (\currule\;$i$) \eTD\eTR
\stopAlethe

\stopRuleDescription

\startRuleDescription{not_xor1}
\startAlethe
\setupTABLE[c][3][width=26em]
\bTR\bTD $i$.\eTD\bTD \ctxsep    \eTD\bTD
$\neg(\lsymb{xor}\,\varphi_1\,\varphi_2)$
\eTD\bTD (\dots) \eTD\eTR
\bTR\bTD $j$.\eTD\bTD \ctxsep    \eTD\bTD
$\varphi_1 , \neg\varphi_2$
\eTD\bTD (\currule\;$i$) \eTD\eTR
\stopAlethe

\noindent
\stopRuleDescription

\startRuleDescription{not_xor2}
\startAlethe
\setupTABLE[c][3][width=26em]
\bTR\bTD $i$.\eTD\bTD \ctxsep    \eTD\bTD
$\neg(\lsymb{xor}\,\varphi_1\,\varphi_2)$
\eTD\bTD (\dots) \eTD\eTR
\bTR\bTD $j$.\eTD\bTD \ctxsep    \eTD\bTD
$\neg\varphi_1 , \varphi_2$
\eTD\bTD (\currule\;$i$) \eTD\eTR
\stopAlethe
\stopRuleDescription

\startRuleDescription{implies}
\startAlethe
\setupTABLE[c][3][width=26em]
\bTR\bTD $i$.\eTD\bTD \ctxsep    \eTD\bTD
$\varphi_1\rightarrow\varphi_2$
\eTD\bTD (\dots) \eTD\eTR
\bTR\bTD $j$.\eTD\bTD \ctxsep    \eTD\bTD
$\neg\varphi_1 , \varphi_2$
\eTD\bTD (\currule\;$i$) \eTD\eTR
\stopAlethe
\stopRuleDescription

\startRuleDescription{not_implies1}
\startAlethe
\setupTABLE[c][3][width=24em]
\bTR\bTD $i$.\eTD\bTD \ctxsep    \eTD\bTD
$\neg (\varphi_1\rightarrow\varphi_2)$
\eTD\bTD (\dots) \eTD\eTR
\bTR\bTD $j$.\eTD\bTD \ctxsep    \eTD\bTD
$\varphi_1$
\eTD\bTD (\currule\;$i$) \eTD\eTR
\stopAlethe
\stopRuleDescription

\startRuleDescription{not_implies2}
\startAlethe
\setupTABLE[c][3][width=24em]
\bTR\bTD $i$.\eTD\bTD \ctxsep    \eTD\bTD
$\neg (\varphi_1\rightarrow\varphi_2)$
\eTD\bTD (\dots) \eTD\eTR
\bTR\bTD $j$.\eTD\bTD \ctxsep    \eTD\bTD
$\neg\varphi_2$
\eTD\bTD (\currule\;$i$) \eTD\eTR
\stopAlethe
\stopRuleDescription

\startRuleDescription{equiv1}
\startAlethe
\setupTABLE[c][3][width=25em]
\bTR\bTD $i$.\eTD\bTD \ctxsep    \eTD\bTD
$\varphi_1≈\varphi_2$
\eTD\bTD (\dots) \eTD\eTR
\bTR\bTD $j$.\eTD\bTD \ctxsep    \eTD\bTD
$\neg\varphi_1 , \varphi_2$
\eTD\bTD (\currule\;$i$) \eTD\eTR
\stopAlethe
\stopRuleDescription

\page
\startRuleDescription{equiv2}
\startAlethe
\setupTABLE[c][3][width=25em]
\bTR\bTD $i$.\eTD\bTD \ctxsep    \eTD\bTD
$\varphi_1≈\varphi_2$
\eTD\bTD (\dots) \eTD\eTR
\bTR\bTD $j$.\eTD\bTD \ctxsep    \eTD\bTD
$\varphi_1, \neg\varphi_2$
\eTD\bTD (\currule\;$i$) \eTD\eTR
\stopAlethe
\stopRuleDescription

\startRuleDescription{not_equiv1}
\startAlethe
\setupTABLE[c][3][width=25em]
\bTR\bTD $i$.\eTD\bTD \ctxsep    \eTD\bTD
$\neg(\varphi_1≈\varphi_2)$
\eTD\bTD (\dots) \eTD\eTR
\bTR\bTD $j$.\eTD\bTD \ctxsep    \eTD\bTD
$\varphi_1 , \varphi_2$
\eTD\bTD (\currule\;$i$) \eTD\eTR
\stopAlethe
\stopRuleDescription


\startRuleDescription{not_equiv2}
\startAlethe
\setupTABLE[c][3][width=25em]
\bTR\bTD $i$.\eTD\bTD \ctxsep    \eTD\bTD
$\neg(\varphi_1≈\varphi_2)$
\eTD\bTD (\dots) \eTD\eTR
\bTR\bTD $j$.\eTD\bTD \ctxsep    \eTD\bTD
$\neg\varphi_1 , \neg\varphi_2$
\eTD\bTD (\currule\;$i$) \eTD\eTR
\stopAlethe
\stopRuleDescription

\startRuleDescription{and_pos}
\startAlethe
\setupTABLE[c][3][width=26em]
\bTR\bTD $i$.\eTD\bTD \ctxsep    \eTD\bTD
$\neg (\varphi_1 \land \cdots \land \varphi_n) , \varphi_k$
\eTD\bTD \currule \eTD\eTR
\stopAlethe

\noindent
with $1 \leq k \leq n$.
\stopRuleDescription

\startRuleDescription{and_neg}
\startAlethe
\setupTABLE[c][3][width=26em]
\bTR\bTD $i$.\eTD\bTD \ctxsep    \eTD\bTD
$(\varphi_1 \land \cdots \land \varphi_n), \neg\varphi_1, \dots , \neg\varphi_n$
\eTD\bTD \currule \eTD\eTR
\stopAlethe

\stopRuleDescription

\startRuleDescription{or_pos}
\startAlethe
\setupTABLE[c][3][width=26.5em]
\bTR\bTD $i$.\eTD\bTD \ctxsep    \eTD\bTD
$\neg (\varphi_1 \lor \cdots \lor \varphi_n) , \varphi_1 , \dots, \varphi_n	$
\eTD\bTD \currule \eTD\eTR
\stopAlethe
\stopRuleDescription

\startRuleDescription{or_neg}
\startAlethe
\setupTABLE[c][3][width=26.5em]
\bTR\bTD $i$.\eTD\bTD \ctxsep    \eTD\bTD
$(\varphi_1 \lor \cdots \lor \varphi_n) , \neg \varphi_k$
\eTD\bTD \currule \eTD\eTR
\stopAlethe

\noindent
with $1 \leq k \leq n$.
\stopRuleDescription

\startRuleDescription{xor_pos1}
\startAlethe
\setupTABLE[c][3][width=26.5em]
\bTR\bTD $i$.\eTD\bTD \ctxsep    \eTD\bTD
$\neg (\varphi_1\,\lsymb{xor}\,\varphi_2) , \varphi_1 , \varphi_2$
\eTD\bTD \currule \eTD\eTR
\stopAlethe
\stopRuleDescription

\page
\startRuleDescription{xor_pos2}
\startAlethe
\setupTABLE[c][3][width=26.5em]
\bTR\bTD $i$.\eTD\bTD \ctxsep    \eTD\bTD
$\neg (\varphi_1\,\lsymb{xor}\,\varphi_2), \neg \varphi_1, \neg \varphi_2$
\eTD\bTD \currule \eTD\eTR
\stopAlethe
\stopRuleDescription

\startRuleDescription{xor_neg1}
\startAlethe
\setupTABLE[c][3][width=26.5em]
\bTR\bTD $i$.\eTD\bTD \ctxsep    \eTD\bTD
$\varphi_1 \,\lsymb{xor}\, \varphi_2, \varphi_1 , \neg \varphi_2$
\eTD\bTD \currule \eTD\eTR
\stopAlethe
\stopRuleDescription

\startRuleDescription{xor_neg2}
\startAlethe
\setupTABLE[c][3][width=26.5em]
\bTR\bTD $i$.\eTD\bTD \ctxsep    \eTD\bTD
$\varphi_1\,\lsymb{xor}\,\varphi_2, \neg \varphi_1 , \varphi_2$
\eTD\bTD \currule \eTD\eTR
\stopAlethe
\stopRuleDescription

\startRuleDescription{implies_pos}
\startAlethe
\setupTABLE[c][3][width=25.5em]
\bTR\bTD $i$.\eTD\bTD \ctxsep    \eTD\bTD
$\neg (\varphi_1 \rightarrow \varphi_2), \neg \varphi_1 , \varphi_2$
\eTD\bTD \currule \eTD\eTR
\stopAlethe
\stopRuleDescription

\startRuleDescription{implies_neg1}
\startAlethe
\setupTABLE[c][3][width=25em]
\bTR\bTD $i$.\eTD\bTD \ctxsep    \eTD\bTD
$\varphi_1 \rightarrow \varphi_2, \varphi_1$
\eTD\bTD \currule \eTD\eTR
\stopAlethe
\stopRuleDescription

\startRuleDescription{implies_neg2}
\startAlethe
\setupTABLE[c][3][width=25em]
\bTR\bTD $i$.\eTD\bTD \ctxsep    \eTD\bTD
$\varphi_1 \rightarrow \varphi_2, \neg \varphi_2$
\eTD\bTD \currule \eTD\eTR
\stopAlethe
\stopRuleDescription

\startRuleDescription{equiv_pos1}
\startAlethe
\setupTABLE[c][3][width=25.5em]
\bTR\bTD $i$.\eTD\bTD \ctxsep    \eTD\bTD
$\neg (\varphi_1 ≈ \varphi_2) , \varphi_1 , \neg \varphi_2$
\eTD\bTD \currule \eTD\eTR
\stopAlethe
\stopRuleDescription

\startRuleDescription{equiv_pos2}
\startAlethe
\setupTABLE[c][3][width=25.5em]
\bTR\bTD $i$.\eTD\bTD \ctxsep    \eTD\bTD
$\neg (\varphi_1 ≈ \varphi_2) , \neg \varphi_1 , \varphi_2$
\eTD\bTD \currule \eTD\eTR
\stopAlethe
\stopRuleDescription

\startRuleDescription{equiv_neg1}
\startAlethe
\setupTABLE[c][3][width=25.5em]
\bTR\bTD $i$.\eTD\bTD \ctxsep    \eTD\bTD
$\varphi_1 ≈ \varphi_2 , \neg \varphi_1 , \neg \varphi_2$
\eTD\bTD \currule \eTD\eTR
\stopAlethe
\stopRuleDescription

\startRuleDescription{equiv_neg2}
\startAlethe
\setupTABLE[c][3][width=25.5em]
\bTR\bTD $i$.\eTD\bTD \ctxsep    \eTD\bTD
$\varphi_1 ≈ \varphi_2 , \varphi_1 , \varphi_2$
\eTD\bTD \currule \eTD\eTR
\stopAlethe
\stopRuleDescription

\startRuleDescription{ite1}
\startAlethe
\setupTABLE[c][3][width=26em]
\bTR\bTD $i$.\eTD\bTD \ctxsep    \eTD\bTD
$\lsymb{ite}\,\varphi_1\,\varphi_2\,\varphi_3$
\eTD\bTD (\dots) \eTD\eTR
\bTR\bTD $j$.\eTD\bTD \ctxsep    \eTD\bTD
$\varphi_1 , \varphi_3$
\eTD\bTD (\currule\;$i$) \eTD\eTR
\stopAlethe
\stopRuleDescription

\startRuleDescription{ite2}
\startAlethe
\setupTABLE[c][3][width=26em]
\bTR\bTD $i$.\eTD\bTD \ctxsep    \eTD\bTD
$\lsymb{ite}\,\varphi_1\,\varphi_2\,\varphi_3$
\eTD\bTD (\dots) \eTD\eTR
\bTR\bTD $j$.\eTD\bTD \ctxsep    \eTD\bTD
$\neg\varphi_1 , \varphi_2$
\eTD\bTD (\currule\;$i$) \eTD\eTR
\stopAlethe
\stopRuleDescription

\startRuleDescription{ite_pos1}
\startAlethe
\setupTABLE[c][3][width=26em]
\bTR\bTD $i$.\eTD\bTD \ctxsep    \eTD\bTD
      $\neg (\lsymb{ite}\,\varphi_1\,\varphi_2\,\varphi_3) , \varphi_1 , \varphi_3$
\eTD\bTD (\currule) \eTD\eTR
\stopAlethe
\stopRuleDescription

\startRuleDescription{ite_pos2}
\startAlethe
\setupTABLE[c][3][width=26em]
\bTR\bTD $i$.\eTD\bTD \ctxsep    \eTD\bTD
      $\neg (\lsymb{ite}\,\varphi_1\,\varphi_2\,\varphi_3) , \neg \varphi_1 , \varphi_2$
\eTD\bTD (\currule) \eTD\eTR
\stopAlethe
\stopRuleDescription

\startRuleDescription{ite_neg1}
\startAlethe
\setupTABLE[c][3][width=26em]
\bTR\bTD $i$.\eTD\bTD \ctxsep    \eTD\bTD
      $\lsymb{ite}\,\varphi_1\,\varphi_2\,\varphi_3 , \varphi_1 , \neg \varphi_3$
\eTD\bTD (\currule) \eTD\eTR
\stopAlethe
\stopRuleDescription

\startRuleDescription{ite_neg2}
\startAlethe
\setupTABLE[c][3][width=26em]
\bTR\bTD $i$.\eTD\bTD \ctxsep    \eTD\bTD
      $\lsymb{ite}\,\varphi_1\,\varphi_2\,\varphi_3 , \neg \varphi_1 , \neg \varphi_2$
\eTD\bTD (\currule) \eTD\eTR
\stopAlethe
\stopRuleDescription

\startRuleDescription{not_ite1}
\startAlethe
\setupTABLE[c][3][width=26em]
\bTR\bTD $i$.\eTD\bTD \ctxsep    \eTD\bTD
$\neg(\lsymb{ite}\,\varphi_1\,\varphi_2\,\varphi_3)$
\eTD\bTD (\dots) \eTD\eTR
\bTR\bTD $j$.\eTD\bTD \ctxsep    \eTD\bTD
$\varphi_1 , \neg\varphi_3$
\eTD\bTD (\currule\;$i$) \eTD\eTR
\stopAlethe
\stopRuleDescription

\startRuleDescription{not_ite2}
\startAlethe
\setupTABLE[c][3][width=26em]
\bTR\bTD $i$.\eTD\bTD \ctxsep    \eTD\bTD
$\neg(\lsymb{ite}\,\varphi_1\,\varphi_2\,\varphi_3)$
\eTD\bTD (\dots) \eTD\eTR
\bTR\bTD $j$.\eTD\bTD \ctxsep    \eTD\bTD
$\neg\varphi_1 , \neg\varphi_2$
\eTD\bTD (\currule\;$i$) \eTD\eTR
\stopAlethe

\stopRuleDescription

\page
\startRuleDescription{connective_def}
  This rule is used to replace connectives by their definition. It can be one
  of the following:

\startAletheS
\setupTABLE[c][4][width=23em]
\bTR\bTD $i$.\eTD\bTD \ctxsep \eTD\bTD $\Gamma$ \eTD\bTD
    $\varphi_1\,\lsymb{xor}\,\varphi_2 ≈
    (\neg\varphi_1 \land \varphi_2) \lor (\varphi_1 \land \neg\varphi_2)$
\eTD\bTD \currule \eTD\eTR
\stopAletheS

\startAletheS
\setupTABLE[c][4][width=23em]
\bTR\bTD $i$.\eTD\bTD \ctxsep \eTD\bTD $\Gamma$ \eTD\bTD
      $\varphi_1≈\varphi_2 ≈
      (\varphi_1 \rightarrow \varphi_2) \land (\varphi_2 \rightarrow \varphi_1)$
\eTD\bTD \currule \eTD\eTR
\stopAletheS

\startAletheS
\setupTABLE[c][4][width=23em]
\bTR\bTD $i$.\eTD\bTD \ctxsep \eTD\bTD $\Gamma$ \eTD\bTD
      $\lsymb{ite}\,\varphi_1\,\varphi_2\,\varphi_3 ≈
      (\varphi_1 \rightarrow \varphi_2) \land (\neg\varphi_1 \rightarrow \varphi_3)$
\eTD\bTD \currule \eTD\eTR
\stopAletheS

\startAletheS
\setupTABLE[c][4][width=23em]
\bTR\bTD $i$.\eTD\bTD \ctxsep \eTD\bTD $\Gamma$ \eTD\bTD
      $\forall x_1, \dots, x_n.\,\varphi ≈ \neg(\exists x_1, \dots, x_n.\,
      \neg\varphi)$
\eTD\bTD \currule \eTD\eTR
\stopAletheS

\stopRuleDescription


\startRuleDescription{and_simplify}
This rule simplifies an $\land$ term by applying equivalent
transformations as long as possible. Hence, the general form is

\startAletheS
\setupTABLE[c][4][width=23em]
\bTR\bTD $i$.\eTD\bTD \ctxsep \eTD\bTD $\Gamma$ \eTD\bTD
$\varphi_1\land \cdots\land\varphi_n ≈ \psi$
\eTD\bTD \currule \eTD\eTR
\stopAletheS

\noindent
where $\psi$ is the transformed term.

The possible transformations are:
\startitemize[packed]
    \item $\top \land \cdots \land \top ⇒ \top$
    \item $\varphi_1 \land \cdots \land \varphi_n ⇒ \varphi_1
    \land \cdots \land \varphi_{n'} $ where the right-hand side has all
    $\top$ literals removed.
    \item $\varphi_1 \land \cdots \land \varphi_n ⇒ \varphi_1
    \land \cdots \land \varphi_{n'} $ where the right-hand side has all
    repeated literals removed.
    \item $\varphi_1 \land\cdots\land \bot\land\cdots \land \varphi_n ⇒ \bot$
    \item $\varphi_1 \land\cdots\land \varphi_i\land \cdots \land \varphi_j\land\cdots \land \varphi_n ⇒ \bot$
  and $\varphi_i$, $\varphi_j$ are such that
  \startformula\startalign
  \NC \varphi_i \NC= \underbrace{\neg \dots \neg}_n \psi \NR
  \NC \varphi_j \NC= \underbrace{\neg \dots \neg}_m \psi
  \stopalign\stopformula
  and one of $n, m$ is odd and the other even.  Either can be $0$.
\stopitemize
\stopRuleDescription

\startRuleDescription{or_simplify}
This rule simplifies an $\lor$ term by applying equivalent
transformations as long as possible. Hence, the general form is

\startAletheS
\setupTABLE[c][4][width=23em]
\bTR\bTD $i$.\eTD\bTD \ctxsep \eTD\bTD $\Gamma$ \eTD\bTD
$(\varphi_1\lor \cdots\lor\varphi_n) ≈ \psi$
\eTD\bTD \currule \eTD\eTR
\stopAletheS

\noindent
where $\psi$ is the transformed term.

The possible transformations are:
\startitemize[packed]
    \item $\bot \lor \cdots \lor \bot ⇒ \bot$
    \item $\varphi_1 \lor \cdots \lor \varphi_n ⇒ \varphi_1
    \lor \cdots \lor \varphi_{n'} $ where the right-hand side has all
    $\bot$ literals removed.
    \item $\varphi_1 \lor \cdots \lor \varphi_n ⇒ \varphi_1
    \lor \cdots \lor \varphi_{n'} $ where the right-hand side has all
    repeated literals removed.
    \item $\varphi_1 \lor\cdots\lor \top\lor\cdots \lor \varphi_n ⇒ \top$
    \item $\varphi_1 \lor\cdots\lor \varphi_i\lor \cdots \lor \varphi_j\lor\cdots \lor \varphi_n ⇒ \top$
  and $\varphi_i$, $\varphi_j$ are such that
  \startformula\startalign
  \NC \varphi_i \NC= \underbrace{\neg \dots \neg}_n \psi \NR
  \NC \varphi_j \NC= \underbrace{\neg \dots \neg}_m \psi
  \stopalign\stopformula
  and one of $n, m$ is odd and the other even.  Either can be $0$.
\stopitemize
\stopRuleDescription

\startRuleDescription{not_simplify}
This rule simplifies an $\neg$ term by applying equivalent
transformations as long as possible. Hence, the general form is

\startAletheS
\setupTABLE[c][4][width=23em]
\bTR\bTD $i$.\eTD\bTD \ctxsep \eTD\bTD $\Gamma$ \eTD\bTD
$\neg\varphi ≈ \psi$
\eTD\bTD \currule \eTD\eTR
\stopAletheS

\noindent
where $\psi$ is the transformed term.

The possible transformations are:
\startitemize[packed]
    \item $\neg (\neg \varphi) ⇒ \varphi$
    \item $\neg \bot ⇒ \top$
    \item $\neg \top ⇒ \bot$
\stopitemize
\stopRuleDescription

\startRuleDescription{implies_simplify}
This rule simplifies an $\rightarrow$ term by applying equivalent
transformations as long as possible. Hence, the general form is

\startAletheS
\setupTABLE[c][4][width=23em]
\bTR\bTD $i$.\eTD\bTD \ctxsep \eTD\bTD $\Gamma$ \eTD\bTD
$\varphi_1\rightarrow \varphi_2 ≈ \psi$
\eTD\bTD \currule \eTD\eTR
\stopAletheS

\noindent
where $\psi$ is the transformed term.

The possible transformations are:
\startitemize[packed]
    \item $\neg \varphi_1 \rightarrow \neg \varphi_2 ⇒  \varphi_2\rightarrow \varphi_1$
    \item $\bot \rightarrow  \varphi ⇒ \top$
    \item $ \varphi \rightarrow \top ⇒ \top$
    \item $\top \rightarrow  \varphi ⇒  \varphi$
		\item $ \varphi \rightarrow \bot ⇒ \neg \varphi$
		\item $ \varphi \rightarrow  \varphi ⇒ \top$
		\item $\neg \varphi \rightarrow  \varphi ⇒  \varphi$
		\item $ \varphi \rightarrow \neg \varphi ⇒ \neg \varphi$
		\item $( \varphi_1\rightarrow \varphi_2)\rightarrow \varphi_2 ⇒  \varphi_1\lor \varphi_2$
\stopitemize
\stopRuleDescription

\startRuleDescription{equiv_simplify}
This rule simplifies a formula with the head symbol $≈\,: \lsymb{Bool}\,\lsymb{Bool}\,\lsymb{Bool}$
by applying equivalent
transformations as long as possible. Hence, the general form is

\startAletheS
\setupTABLE[c][4][width=23em]
\bTR\bTD $i$.\eTD\bTD \ctxsep \eTD\bTD $\Gamma$ \eTD\bTD
$(\varphi_1≈ \varphi_2) ≈ \psi$
\eTD\bTD \currule \eTD\eTR
\stopAletheS

\noindent
where $\psi$ is the transformed term.

The possible transformations are:
\startitemize[packed]
    \item $(\neg \varphi_1 ≈ \neg \varphi_2) ⇒ ( \varphi_1≈ \varphi_2)$
    \item $( \varphi≈ \varphi) ⇒ \top$
    \item $( \varphi≈ \neg \varphi) ⇒ \bot$
    \item $(\neg \varphi≈  \varphi) ⇒ \bot$
    \item $(\top ≈  \varphi) ⇒  \varphi$
    \item $( \varphi ≈ \top) ⇒  \varphi$
    \item $(\bot ≈  \varphi) ⇒ \neg \varphi$
    \item $( \varphi ≈ \bot) ⇒ \neg \varphi$
\stopitemize
\stopRuleDescription

\startRuleDescription{bool_simplify}
This rule simplifies a boolean term by applying equivalent
transformations as long as possible. Hence, the general form is

\startAletheS
\setupTABLE[c][4][width=23em]
\bTR\bTD $i$.\eTD\bTD \ctxsep \eTD\bTD $\Gamma$ \eTD\bTD
$\varphi≈ \psi$
\eTD\bTD \currule \eTD\eTR
\stopAletheS

\noindent
where $\psi$ is the transformed term.

The possible transformations are:
\startitemize[packed]
	\item $\neg(\varphi_1\rightarrow \varphi_2) ⇒ (\varphi_1 \land \neg \varphi_2)$
	\item $\neg(\varphi_1\lor \varphi_2) ⇒ (\neg \varphi_1 \land \neg \varphi_2)$
	\item $\neg(\varphi_1\land \varphi_2) ⇒ (\neg \varphi_1 \lor \neg \varphi_2)$
	\item $(\varphi_1 \rightarrow (\varphi_2\rightarrow \varphi_3)) ⇒ (\varphi_1\land \varphi_2) \rightarrow \varphi_3$
	\item $((\varphi_1\rightarrow \varphi_2)\rightarrow \varphi_2)  ⇒ (\varphi_1\lor \varphi_2)$
	\item $(\varphi_1 \land (\varphi_1\rightarrow \varphi_2)) ⇒ (\varphi_1 \land \varphi_2)$
	\item $((\varphi_1\rightarrow \varphi_2) \land \varphi_1) ⇒ (\varphi_1 \land \varphi_2)$
\stopitemize
\stopRuleDescription

\page
\startRuleDescription{ac_simp}
  This rule simplifies nested occurrences of $\lor$ or $\land$:

\startAletheS
\setupTABLE[c][4][width=23em]
\bTR\bTD $i$.\eTD\bTD \ctxsep \eTD\bTD $\Gamma$ \eTD\bTD
      $\psi  ≈ \varphi_1 \circ\cdots\circ\varphi_n$
\eTD\bTD \currule \eTD\eTR
\stopAletheS

\noindent
  where $\circ \in \{\lor, \land\}$ and $\psi$ is a nested application of $\circ$.
  The literals $\varphi_i$ are literals of the flattening of $\psi$ with duplicates
  removed.
\stopRuleDescription

\startRuleDescription{ite_simplify}
  This rule simplifies an if-then-else term by applying equivalent
  transformations until fix point\footnote{Note however that the order of the
    application is important, since the set of rules is not confluent. For
    example, the term $(\lsymb{ite} \top \; t_1 \; t_2 ≈
    t_1)$ can be simplified into both $p$ and $(\neg (\neg p))$ depending on the
    order of applications.}
  %
It has the form
\startAletheS
\setupTABLE[c][4][width=23em]
\bTR\bTD $i$.\eTD\bTD \ctxsep \eTD\bTD $\Gamma$ \eTD\bTD
$\lsymb{ite}\,\varphi\,t_1\,t_2 ≈ u$
\eTD\bTD \currule \eTD\eTR
\stopAletheS

\noindent
where $u$ is the transformed term.

The possible transformations are:
\startitemize[packed]
    \item $\lsymb{ite}\, \top      \, t_1 \, t_2 ⇒ t_1$
    \item $\lsymb{ite}\, \bot      \, t_1 \, t_2 ⇒ t_2$
    \item $\lsymb{ite}\, \psi      \, t \, t ⇒ t$
    \item $\lsymb{ite}\, \neg \varphi \, t_1 \, t_2 ⇒ \lsymb{ite}\, \varphi \, t_2 \, t_1$
    \item $\lsymb{ite}\, \psi \, (\lsymb{ite}\, \psi\,t_1\,t_2)\, t_3 ⇒
			\lsymb{ite}\, \psi\, t_1\, t_3$
    \item $\lsymb{ite}\, \psi \, t_1\, (\lsymb{ite}\, \psi\,t_2\,t_3) ⇒
			\lsymb{ite}\, \psi\, t_1\, t_3$
    \item $\lsymb{ite}\, \psi \, \top\, \bot ⇒ \psi$
    \item $\lsymb{ite}\, \psi \, \bot\, \top ⇒ \neg\psi$
    \item $\lsymb{ite}\, \psi \, \top \, \varphi ⇒ \psi\lor\varphi$
    \item $\lsymb{ite}\, \psi \, \varphi\,\bot ⇒ \psi\land\varphi$
    \item $\lsymb{ite}\, \psi \, \bot\, \varphi ⇒ \neg\psi\land\varphi$
    \item $\lsymb{ite}\, \psi \, \varphi\,\top ⇒ \neg\psi\lor\varphi$
\stopitemize
\stopRuleDescription

\startRuleDescription{qnt_simplify}
  This rule simplifies a $\forall$-formula with a constant predicate.

\startAletheS
\setupTABLE[c][4][width=23em]
\bTR\bTD $i$.\eTD\bTD \ctxsep \eTD\bTD $\Gamma$ \eTD\bTD
      $\forall x_1, \dots, x_n. \varphi ≈ \varphi$
\eTD\bTD \currule \eTD\eTR
\stopAletheS

\noindent
  where $\varphi$ is either $\top$ or $\bot$.
\stopRuleDescription

\page
\startRuleDescription{onepoint}
The {\currule} rule is the \quotation{one-point-rule}. That is: it eliminates quantified
variables that can only have one value.

\startAletheS
\setupTABLE[c][4][width=10em]
\aletheLineS
\bTR[bottomframe=on] \bTD[bottomframe=off] $j$.\eTD\bTD[leftframe=on]
  $\Gamma, x_{k_1},\dots, x_{k_m},  x_{j_1} \mapsto t_{j_1}, \dots ,  x_{j_o} \mapsto t_{j_o}$
	\eTD\bTD \ctxsep
  \eTD\bTD
  $\varphi ≈ \varphi'$
  \eTD\bTD (\dots) \eTD\eTR
\bTR\bTD $k$.\eTD \bTD \eTD
    \bTD \ctxsep \eTD\bTD
$Q x_1, \dots, x_n.\varphi ≈ Q x_{k_1}, \dots, x_{k_m}. \varphi'$
    \eTD\bTD \currule{} \eTD\eTR
     \eTD\bTD \eTD\eTR
\stopAletheS

\noindent
where $Q\in\{\forall, \exists\}$,  $n = m + o$,  $k_1, \dots, k_m$ and
$j_1, \dots, j_o$ are monotone
mappings to $1, \dots, n$, and no $x_{k_i}$ appears in $x_{j_1}, \dots, x_{j_o}$.

The terms $t_{j_1}, \dots, t_{j_o}$ are the points of the variables
$x_{j_1}, \dots, x_{j_o}$. Points are defined by equalities $x_i≈ t_i$
with positive polarity in the term $\varphi$.

\tightpara{Remark.} Since an eliminated variable $x_i$ might appear free in a
term $t_j$, it is necessary to replace $x_i$ with $t_i$ inside $t_j$. While
this substitution is performed correctly, the proof for it is currently
missing.

\stopRuleDescription

\startruleexample
An application of the \currule{} rule on the term $(\forall x, y.\, x ≈ y
\rightarrow (f\,x)\land (f\,y))$ look like this:

\startSMTCode
(anchor :step t3 :args ((:= y x)))
(step t3.t1 (cl (= x y)) :rule refl)
(step t3.t2 (cl (= (= x y) (= x x)))
    :rule cong :premises (t3.t1))
(step t3.t3 (cl (= x y)) :rule refl)
(step t3.t4 (cl (= (f y) (f x)))
    :rule cong :premises (t3.t3))
(step t3.t5 (cl (= (and (f x) (f y)) (and (f x) (f x))))
    :rule cong :premises (t3.t4))
(step t3.t6 (cl (= (=> (= x y) (and (f x) (f y)))
                   (=> (= x x) (and (f x) (f x)))))
    :rule cong :premises (t3.t2 t3.t5))
(step t3 (cl (=
        (forall ((x S) (y S)) (=> (= x y) (and (f x) (f y))))
        (forall ((x S))       (=> (= x x) (and (f x) (f x))))))
    :rule qnt_simplify)
\stopSMTCode
\stopruleexample

\page
\startRuleDescription{qnt_join}
\startAletheS
\setupTABLE[c][4][width=26em]
\bTR\bTD $i$.\eTD\bTD \ctxsep \eTD\bTD $\Gamma$ \eTD\bTD
      $Q x_1, \dots, x_n.\,(Q x_{n+1}, \dots, x_{m}.\,\varphi)
      ≈
      Q x_{k_1}, \dots, x_{k_o}.\,\varphi$
\eTD\bTD \currule \eTD\eTR
\stopAletheS

\noindent
  where $m > n$ and $Q\in\{\forall, \exists\}$. Furthermore, $k_1, \dots, k_o$ is a monotonic
  map to $1, \dots, m$ such that $x_{k_1}, \dots, x_{k_o}$ are pairwise
  distinct, and $\{x_1, \dots, x_m\} = \{x_{k_1}, \dots, x_{k_o}\}$.
\stopRuleDescription

\startRuleDescription{qnt_rm_unused}
\startAletheS
\setupTABLE[c][4][width=23em]
\bTR\bTD $i$.\eTD\bTD \ctxsep \eTD\bTD $\Gamma$ \eTD\bTD
      $Q x_1, \dots, x_n.\,\varphi ≈ Q x_{k_1}, \dots, x_{k_m}.\,\varphi$
\eTD\bTD \currule \eTD\eTR
\stopAletheS

  \noindent
  where $m \leq n$ and $Q\in\{\forall, \exists\}$. Furthermore, $k_1, \dots, k_m$ is
  a monotonic map to $1, \dots, n$ and if $x\in \{x_j\; |\; j \in \{1, \dots,
  n\} \land j\in\not \{k_1, \dots, k_m\}\}$ then $x$ is not free in $P$.

\stopRuleDescription

\startRuleDescription{eq_simplify}
  This rule simplifies an $≈$ term by applying equivalent
  transformations as long as possible. Hence, the general form is

\startAletheS
\setupTABLE[c][4][width=25em]
\bTR\bTD $i$.\eTD\bTD \ctxsep \eTD\bTD $\Gamma$ \eTD\bTD
 $(t_1≈ t_2) ≈ \varphi$
\eTD\bTD \currule \eTD\eTR
\stopAletheS

\noindent
  where $\psi$ is the transformed term.

  The possible transformations are:
  \startitemize[packed]
  \item $t ≈ t ⇒ \top$
  \item $(t_1 ≈ t_2) ⇒ \bot$ if $t_1$ and $t_2$ are different numeric constants.
  \item $\neg (t ≈ t) ⇒ \bot$ if $t$ is a numeric constant.
  \stopitemize
\stopRuleDescription

\startRuleDescription{div_simplify}
This rule simplifies a division by applying equivalent
transformations. The general form is

\startAletheS
\setupTABLE[c][4][width=25em]
\bTR\bTD $i$.\eTD\bTD \ctxsep \eTD\bTD $\Gamma$ \eTD\bTD
$(t_1\, /\,  t_2) ⇒ t_3$
\eTD\bTD \currule \eTD\eTR
\stopAletheS

\noindent The possible transformations are:
\startitemize[packed]
  \item $t\, /\, t ⇒ 1$
  \item $t\, /\, 1 ⇒ t$
	\item $t_1\,  /\,  t_2 ⇒ t_3$
		if $t_1$ and $t_2$ are constants and $t_3$ is $t_1$
		divided by $t_2$ according to the semantic of the current theory.
\stopitemize
\stopRuleDescription

\page
\startRuleDescription{prod_simplify}
This rule simplifies a product by applying equivalent
transformations as long as possible. The general form is

\startAletheS
\setupTABLE[c][4][width=25em]
\bTR\bTD $i$.\eTD\bTD \ctxsep \eTD\bTD $\Gamma$ \eTD\bTD
$t_1\times\cdots\times t_n ≈ u$
\eTD\bTD \currule \eTD\eTR
\stopAletheS

\noindent
where $u$ is either a constant or a product.

The possible transformations are:
\startitemize[packed]
    \item $t_1\times\cdots\times t_n ⇒ u$ where all
    $t_i$ are constants and $u$ is their product.
    \item $t_1\times\cdots\times t_n ⇒ 0$ if any
    $t_i$ is $0$.
    \item $t_1\times\cdots\times t_n ⇒
			c \times t_{k_1}\times\cdots\times t_{k_n}$ where $c$
			is the product of the constants of $t_1, \dots, t_n$ and
			$t_{k_1}, \dots, t_{k_n}$ is $t_1, \dots, t_n$
			with the constants removed.
    \item $t_1\times\cdots\times t_n ⇒
			t_{k_1}\times\cdots\times t_{k_n}$: same as above if $c$ is $1$.
\stopitemize
\stopRuleDescription

\startRuleDescription{unary_minus_simplify}
This rule is either

\startAletheS
\setupTABLE[c][4][width=20em]
\bTR\bTD $i$.\eTD\bTD \ctxsep \eTD\bTD $\Gamma$ \eTD\bTD
$- (-t) ≈ t$
\eTD\bTD \currule \eTD\eTR
\stopAletheS

\noindent
or
\startAletheS
\setupTABLE[c][4][width=20em]
\bTR\bTD $i$.\eTD\bTD \ctxsep \eTD\bTD $\Gamma$ \eTD\bTD
$-t ≈ u$
\eTD\bTD \currule \eTD\eTR
\stopAletheS

\noindent
where $u$ is the negated numerical constant $t$.
\stopRuleDescription


\startRuleDescription{minus_simplify}
This rule simplifies a subtraction by applying equivalent
transformations. The general form is

\startAletheS
\setupTABLE[c][4][width=22em]
\bTR\bTD $i$.\eTD\bTD \ctxsep \eTD\bTD $\Gamma$ \eTD\bTD
$t_1- t_2 ≈ u$
\eTD\bTD \currule \eTD\eTR
\stopAletheS

\noindent The possible transformations are:
\startitemize[packed]
    \item $t - t ⇒ 0$
    \item $t_1 - t_2 ⇒ t_3$ where $t_1$
    and $t_2$ are numerical constants and $t_3$ is $t_2$ subtracted
    from~$t_1$.
    \item $t - 0 ⇒ t$
    \item $0 - t ⇒ -t$
\stopitemize
\stopRuleDescription

\startRuleDescription{sum_simplify}
This rule simplifies a sum by applying equivalent
transformations as long as possible. The general form is

\startAletheS
\setupTABLE[c][4][width=22em]
\bTR\bTD $i$.\eTD\bTD \ctxsep \eTD\bTD $\Gamma$ \eTD\bTD
$t_1+\cdots+t_n ≈ u$
\eTD\bTD \currule \eTD\eTR
\stopAletheS

\noindent
where $u$ is either a constant or a product.

The possible transformations are:
\startitemize[packed]
    \item $t_1+\cdots+t_n ⇒ c$ where all
    $t_i$ are constants and $c$ is their sum.
    \item $t_1+\cdots+t_n ⇒
			c + t_{k_1}+\cdots+t_{k_n}$ where $c$
			is the sum of the constants of $t_1, \dots, t_n$ and
			$t_{k_1}, \dots, t_{k_n}$ is $t_1, \dots, t_n$
			with the constants removed.
    \item $t_1+\cdots+t_n ⇒
			t_{k_1}+\cdots+t_{k_n}$: same as above if $c$ is
			$0$.
\stopitemize
\stopRuleDescription


\startRuleDescription{comp_simplify}
This rule simplifies a comparison by applying equivalent
transformations as long as possible. The general form is

\startAletheS
\setupTABLE[c][4][width=22em]
\bTR\bTD $i$.\eTD\bTD \ctxsep \eTD\bTD $\Gamma$ \eTD\bTD
$t_1 \bowtie t_n ≈ \psi$
\eTD\bTD \currule \eTD\eTR
\stopAletheS

\noindent
where $\bowtie$ is as for the proof rule \proofRule{la_generic}, but never
$≈$.

The possible transformations are:
\startitemize[packed]
    \item $t_1 < t_2 ⇒ \varphi$ where $t_1$ and
    $t_2$ are numerical constants and $\varphi$ is $\top$ if $t_1$ is
    strictly greater than $t_2$ and $\bot$ otherwise.
    \item $t < t ⇒ \bot$
    \item $t_1 \leq t_2 ⇒ \varphi$ where $t_1$ and
    $t_2$ are numerical constants and $\varphi$ is $\top$ if $t_1$ is
    greater than $t_2$ or equal and $\bot$ otherwise.
    \item $t \leq t ⇒ \top$
    \item $t_1 \geq t_2 ⇒ t_2 \leq t_1$
    \item $t_1 < t_2 ⇒ \neg (t_2 \leq t_1)$
    \item $t_1 > t_2 ⇒ \neg (t_1 \leq t_2)$
\stopitemize
\stopRuleDescription

\startRuleDescription{let}
  This rule eliminates $\lsymb{let}$. It has the form

\startAletheS
\setupTABLE[c][4][width=14.5em]
\bTR \bTD $i_1$.\eTD\bTD
  $\Gamma$
	\eTD\bTD \ctxsep \eTD\bTD
  $t_{1} ≈ s_{1}$
  \eTD\bTD (\dots) \eTD\eTR
\aletheLineSB
\bTR \bTD $i_n$.\eTD\bTD
  $\Gamma$
	\eTD\bTD \ctxsep \eTD\bTD
  $t_{n} ≈ s_{n}$
  \eTD\bTD (\dots) \eTD\eTR
\aletheLineS
\bTR[bottomframe=on] \bTD[bottomframe=off] $j$.\eTD\bTD[leftframe=on]
  $\Gamma, x_1 \mapsto s_1, \dots,  x_n \mapsto s_n$
	\eTD\bTD \ctxsep
  \eTD\bTD
  $u ≈ u'$
  \eTD\bTD (\dots) \eTD\eTR
\bTR\bTD $k$.\eTD \bTD $\Gamma$ \eTD
    \bTD \ctxsep \eTD\bTD
     $(\lsymb{let}\,x_1 = t_1,\, \dots,\, x_n = t_n\,\lsymb{in}\, u) ≈ u'$
    \eTD\bTD (\currule{}\;$i_1$, \dots, $i_n$) \eTD\eTR
     \eTD\bTD \eTD\eTR
\stopAletheS

	The premise $i_1, \dots, i_n$ must be in the same subproof as
	the \currule{} step.  If for $t_i≈ s_i$ the $t_i$ and $s_i$
	are syntactically equal, the premise
  is skipped.
\stopRuleDescription

\page
\startRuleDescription{distinct_elim}
This rule eliminates the $\lsymb{distinct}$ predicate. If called with one
argument this predicate always holds:
\startAletheS
\setupTABLE[c][4][width=22em]
\bTR\bTD $i$.\eTD\bTD \ctxsep \eTD\bTD $\Gamma$ \eTD\bTD
$(\lsymb{distinct}\, t) ≈ \top$
\eTD\bTD \currule \eTD\eTR
\stopAletheS

If applied to terms of type $\lsymb{Bool}$ more than two terms can never be
distinct, hence only two cases are possible:

\startAletheS
\setupTABLE[c][4][width=22em]
\bTR\bTD $i$.\eTD\bTD \ctxsep \eTD\bTD $\Gamma$ \eTD\bTD
$(\lsymb{distinct}\,\varphi\,\psi) ≈ \neg (\varphi ≈ \psi)$
\eTD\bTD \currule \eTD\eTR
\stopAletheS

\noindent
and
\startAletheS
\setupTABLE[c][4][width=22em]
\bTR\bTD $i$.\eTD\bTD \ctxsep \eTD\bTD $\Gamma$ \eTD\bTD
$(\lsymb{distinct}\,\varphi_1\,\varphi_2\,\varphi_3\,\dots) ≈ \bot$
\eTD\bTD \currule \eTD\eTR
\stopAletheS

The general case is
\startAletheS
\setupTABLE[c][4][width=22em]
\bTR\bTD $i$.\eTD\bTD \ctxsep \eTD\bTD $\Gamma$ \eTD\bTD
$(\lsymb{distinct}\,t_1\,\dots\, t_n) ≈
\bigwedge_{i=1}^{n}\bigwedge_{j=i+1}^{n} t_i\;{≈\not}\;t_j$
\eTD\bTD \currule \eTD\eTR
\stopAletheS
\stopRuleDescription

\startRuleDescription{la_rw_eq}
\startAlethe
\setupTABLE[c][3][width=26em]
\bTR\bTD $i$.\eTD\bTD \ctxsep  \eTD\bTD
$(t ≈ u) ≈ (t \leq u \land u \leq t)$
\eTD\bTD \currule \eTD\eTR
\stopAlethe

\tightpara{Remark.} While the connective could be $≈$,
currently an equality is used.
\stopRuleDescription

\startRuleDescription{nary_elim}
This rule replaces $n$-ary operators with their equivalent
application of the binary operator. It is never applied to $\land$ or $\lor$.

Three cases are possible.
If the operator $\circ$ is left associative, then the rule has the form
\startAletheS
\setupTABLE[c][4][width=25em]
\bTR\bTD $i$.\eTD\bTD \ctxsep \eTD\bTD $\Gamma$ \eTD\bTD
$\bigcirc_{i=1}^{n} t_i ≈
(\dots( t_1\circ  t_2) \circ  t_3)\circ \cdots  t_n)$
\eTD\bTD \currule \eTD\eTR
\stopAletheS

If the operator $\circ$ is right associative, then the rule has the form
\startAletheS
\setupTABLE[c][4][width=25em]
\bTR\bTD $i$.\eTD\bTD \ctxsep \eTD\bTD $\Gamma$ \eTD\bTD
$\bigcirc_{i=1}^{n} t_i ≈
( t_1 \circ \cdots \circ ( t_{n-2} \circ ( t_{n-1} \circ  t_n)\dots)$
\eTD\bTD \currule \eTD\eTR
\stopAletheS

If the operator is {\em chainable}, then it has the form

\startAletheS
\setupTABLE[c][4][width=25em]
\bTR\bTD $i$.\eTD\bTD \ctxsep \eTD\bTD $\Gamma$ \eTD\bTD
$\bigcirc_{i=1}^{n} t_i ≈
( t_1\circ t_2) \land ( t_2 \circ  t_3) \land \cdots
\land ( t_{n-1}\circ t_n)$
\eTD\bTD \currule \eTD\eTR
\stopAletheS
\stopRuleDescription

\startRuleDescription{bfun_elim}
\startAlethe
\setupTABLE[c][3][width=24em]
\bTR\bTD $i$.\eTD\bTD \ctxsep  \eTD\bTD
$\psi$
\eTD\bTD (\dots) \eTD\eTR
\bTR\bTD $j$.\eTD\bTD \ctxsep  \eTD\bTD
$\varphi$
\eTD\bTD (\currule\; $i$) \eTD\eTR
\stopAlethe

The formula $\varphi$ is $\psi$ after boolean functions have been simplified.
This happens in a two step process. Both steps recursively iterate over $\psi$.
The first step expands quantified variable of type $\lsymb{Bool}$. Hence,
$(\exists x.\,t)$ becomes $t[x\mapsto \bot]\lor t[x\mapsto \top]$ and
$(\forall x.\,t)$ becomes $t[x\mapsto \bot]\land t[x\mapsto \top]$. If $n$ variables of sort
$\lsymb{Bool}$ appear in a quantifier, the disjunction (conjunction) has
$2^n$ terms. Each term replaces the variables in $t$ according
to the bits of a number which is increased by one for each subsequent
term starting from zero. The left-most variable corresponds to the
least significant bit.

The second step expands function argument of boolean types by introducing
appropriate if-then-else terms. For example, consider $(f\,x\, P\, y)$ where
$P$ is some formula. Then we replace this term by $(\lsymb{ite}\, P\,
(f\,x\, \top\ y)\,(f\,x\, \bot\, y))$. If the argument is already the constant $\top$
or $\bot$, it is ignored.

\stopRuleDescription

\startRuleDescription{ite_intro}
\startAlethe
\setupTABLE[c][3][width=24em]
\bTR\bTD $i$.\eTD\bTD \ctxsep  \eTD\bTD
$t ≈ (t' \land u_1 \land \dots \land u_n)$
\eTD\bTD (\currule) \eTD\eTR
\stopAlethe

The term $t$ (the formula $\varphi$) contains the $\lsymb{ite}$ operator.
Let $s_1, \dots, s_n$ be the terms starting with $\lsymb{ite}$, i.e.
$s_i := \lsymb{ite}\,\psi_i\,r_i\,r'_i$, then $u_i$ has the form
\startformula
	\lsymb{ite}\,\psi_i\,(s_i ≈ r_i)\,(s_i ≈ r'_i)
\stopformula

\noindent
The term $t'$ is equal to the term $t$ up to the
reordering of equalities where one argument is an $\lsymb{ite}$
term.

\tightpara{Remark.} This rule stems from the introduction of fresh
constants for if-then-else terms inside veriT. Internally $s_i$ is a new
constant symbol and the $\varphi$ on the right side of the equality is
$\varphi$ with the if-then-else terms replaced by the constants. Those
constants are unfolded during proof printing. Hence, the slightly strange
form and the reordering of equalities.
\stopRuleDescription

\page % TODO: double sided
\subject[sec:alethe:rules-index]{Rule Index}
\placeruleIndex[textstyle={\ss\addff{table}},style={\bf\WORDS}]

\stoptitle
\stopcomponent

